\documentclass[10pt,a4paper]{article}
\usepackage[top=60pt,bottom=60pt,left=60pt,right=60pt]{geometry}
\usepackage{amsmath}
%\usepackage{draftwatermark}
\usepackage{xcolor}



%\SetWatermarkText{SOLUTIONS}
%\SetWatermarkScale{4.0} % Adjust size (default is 1)
%\SetWatermarkAngle{45} % Diagonal angle in degrees (default is 45)
% \SetWatermarkLightness{0.8} % Gray level from 0 (black) to 1 (white)
%\SetWatermarkColor{red!20} % Set color to red with 20% opacity (transparent red)



\begin{document}

\noindent 2026.11.20-22 \hfill \textbf{PushQuantum 2026 Hackathon -- Hints} \hfill Aqarios GmbH

\begin{center}
    \Large See further hints ($XY$-mixers, warm-starts, etc.) in code files.
\end{center}

\section*{\texttt{qaoa.py}}


\begin{itemize}
    \item[1.] \textbf{QAOA Parameter Initialization.}
    Take a closer look at the function
    \texttt{\_initialize\_start\_parameters}.
    Are the current parameter ranges appropriate?
    Consider the periodicity of the corresponding
    operators and whether it allows you to reduce
    the search space. Keep in mind that different
    mixer types may have different periodicities.



    \item[2.] \textbf{Explore Alternative Mixer Layers.}
    Certain optimization problems impose constraints
    on the number of active bits in a solution.
    For example, a feasible bitstring may be
    required to contain exactly $k$ ones.

    Instead of using the standard $X$-mixer,
    consider an \textbf{XY-mixer}, which can
    preserve the Hamming weight of a quantum state
    and therefore keep the evolution within the
    feasible subspace.

    You can find a relevant paper in the
    \texttt{Papers} folder.



    \item[3.] \textbf{Optimize Fewer Parameters.}
    You have already encountered different
    parameter initialization strategies, such as
    random initialization and ramp schedules.
    But what if you could reduce the number of
    parameters that need to be optimized?

    One approach is the \textbf{Linear Ramp}
    parameterization, along with its many
    possible variations.

    The key idea is to optimize only a small
    set of parameters and derive the remaining
    QAOA angles from them. This reduces the
    dimensionality of the optimization problem
    and may significantly lower the computational
    cost of classical optimization.

    You can find more information about Linear
    Ramp in the \texttt{Papers} folder, as well
    as numerous related approaches on arXiv.

\end{itemize}




\section*{\texttt{custom\_vqe.py}}

\begin{itemize}
    \item[1.] \texttt{\_construct\_hardware\_efficient}: Do all gates have an effect on the probabilities? Unnecessary gates may increase the runtime.



    \item[2.] \texttt{\_construct\_real\_amplitudes}: Can the $CX$ block be simplified?
    

    
    \item[3.] Expectation matters: Are there other ways such as CVAR and Mean to implement \texttt{cost\_function}?
    


\end{itemize}


\section*{Additional ideas}

\begin{itemize}
    \item[1.] Think about the scaling of the model (penalty terms, overall scaling). A similar scaling of the objective over all problems can have some advantages.
    \item[2.] How does the performance depend on the number of layers?
    \item[3.] There are not only $XY$-mixers, e.g.\ Dicke mixers, 2-way $XY$-mixers, etc. There is also MOQAOA: constraints without slack, penalty is tuned beforehand. You may also consider counter-diabatic mixers.
    \item[4.] You may consider not only linear ramps, but also non-linear ramps. 
\end{itemize}



\end{document}
